\subsection{Comparison with CASAAL: the whole translation}
\label{sec:casaal-eval}
CASAAL~\cite{casaal_tool} implements the MTL$_{0,\infty}$ construction of
Li et al.~\cite{li_spin7}. It is not an optimizer that could be run on our
automata: it builds its own automata, so it is a reference for the whole
translation, from the requirement to the automaton given to a timed model
checker, and not for the optimization passes alone. We therefore compare the
two tools on four measures: the size of the automaton, the time to generate
it, the number of states a timed model checker stores on it, and, to say what
the optimizer contributes, the same measures with the optimizer disabled
(\texttt{-n 0}). We ran CASAAL with mutually exclusive events, which matches
the one-event-per-position words of our semantics. Its acceptance is carried
by transitions and ours by locations; sizes are compared after converting its
automata to location acceptance (Section~\ref{sec:casaal-size}).

Two configurations of our tool appear in the tables. \emph{Auto} is the
default, which runs Spot at level \texttt{low} for at most 5~s and then falls
back to level \texttt{any} on the simplified formula with Spot's simulation
reductions switched off (Section~\ref{sec:spot-level}). \emph{Manual} is the
explicit setting \texttt{-simp -spot-level any -spot-x ba-simul=0} or
\texttt{-spot-level low}, whichever was faster for that requirement, shown only
when it beats auto by at least 10\% and 0.1~s. The two configurations produce
automata of the same size on every pattern of the tables; they differ in time
only.

\subsubsection{Size}
\label{sec:casaal-size}
\emph{Requirements as written (positive form).} On the 27 patterns that CASAAL
encodes exactly and both tools completed, the totals are 222 locations and 976
transitions for ours (with Spot level \texttt{small}, the former default, which
produced the archived runs) against 225 and 983 for CASAAL after conversion to
location acceptance. Ours is strictly smaller on four patterns (F8,
F11, F12, F14), equal on 21, and larger on two (F5, F6: 7 locations and 25
transitions against 6 and 22). With the current default the same two patterns
have 27 transitions, for a total of 980 against 983. CASAAL is
smaller on the lower-bounded responses F5 and F6, and ours on F8, F11, F12 and
F14, where the optimizer merges locations that CASAAL keeps.

\emph{Negated requirements (what the model checker receives).} Table~\ref{tab:casaal-gen}
gives the sizes for R2--R8 and N5--N8: ours equals CASAAL's on every one
(equal location and transition counts; we did not check isomorphism), and the
optimizer's output is exactly the size of CASAAL's on the nested family
(N5: 461 transitions, N6: 1,541, N7: 5,105). On the 20 patterns whose
products TChecker completes, ours is smaller on F5, F6 (3 locations and 5
transitions against 4 and 7) and F8 (2 and 3 against 4 and 7) and equal on the
others.

\subsubsection{Generation time}
\label{sec:casaal-time}
\begin{table}[tp]
\centering
\caption{Negated requirements: CASAAL against \texttt{mtl2tba}. Size is reachable locations/transitions of the automaton given to the model checker (ours: after optimization). Times are wall-clock seconds for the whole translation (generation plus, for ours, optimization and export), CASAAL on native Windows and \texttt{mtl2tba} under WSL~2 on the same machine; CASAAL times are the best of three runs and ours single runs (article Table~\ref{tab:same-driver} gives medians of five runs with one driver). Sizes of CASAAL's automata are after conversion to location acceptance. Spot level: the current default (auto, with the fallback marked) and the manual setting. ``Auto'' is the default of the tool; $^{\dagger}$ marks runs where Spot had not finished within the 5~s budget and the tool fell back to \texttt{any} on the simplified formula. The last column gives the time of the fastest explicit setting when it is at least 10\% and 0.1~s faster than auto (``--'': no explicit setting qualifies, e.g. on N8 the manual run took 243~s against 247.8~s for auto in this table's run): \texttt{low}, or \texttt{-simp -spot-level any -spot-x ba-simul=0}, written \texttt{any}; the sizes of the automata are the same for both settings except where noted in the text.}
\label{tab:casaal-gen}
\scriptsize
\setlength{\tabcolsep}{4pt}
\begin{tabular}{@{}lrrrrr@{}}
\toprule
 & CASAAL size & Ours size & CASAAL time & Ours, auto & Ours, manual (setting)\\
\midrule
R2 & 4/8 & 4/8 & 0.03 & 0.1 & -- \\
R3 & 5/11 & 5/11 & 0.03 & 0.1 & -- \\
R4 & 6/14 & 6/14 & 0.04 & 0.1 & -- \\
R5 & 7/17 & 7/17 & 0.06 & 0.1 & -- \\
R6 & 8/20 & 8/20 & 0.09 & 0.1 & -- \\
R7 & 9/23 & 9/23 & 0.1 & 0.1 & -- \\
R8 & 10/26 & 10/26 & 0.2 & 0.1 & -- \\
N5 & 33/461 & 33/461 & 0.3 & 0.8 & 0.2 (\texttt{any}) \\
N6 & 65/1541 & 65/1541 & 1.9 & 6.6$^{\dagger}$ & 1.3 (\texttt{any}) \\
N7 & 129/5105 & 129/5105 & 11.3 & 22.1$^{\dagger}$ & 15.1 (\texttt{any}) \\
N8 & 257/16769 & 257/16769 & 73.8 & 247.8$^{\dagger}$ & -- \\
\bottomrule
\end{tabular}
\end{table}

Table~\ref{tab:casaal-gen} reports the wall-clock time of the whole
translation of the negated requirement. CASAAL runs natively on Windows and
\texttt{mtl2tba} under WSL~2 on the same machine, so the comparison carries the
overhead of the virtual machine for our tool; the CASAAL times are the best of
three runs, ours single runs.

\emph{Small requirements and conjunctions.} Every negated F and R pattern takes our
tool about 0.1~s, of which almost all is the start of the processes (Spot and
the tool itself): a bounded response takes CASAAL 0.02~s, so on the small
patterns CASAAL is faster in absolute terms, by about 0.1~s. On the
conjunctions of bounded responses CASAAL's time grows with the number of
conjuncts (0.09~s for six responses, 0.13~s for seven, 0.20~s for eight) while
ours stays at about 0.1~s with Spot level \texttt{low}, so that ours is ahead
by 0.02~s at seven and 0.08~s at eight, differences within the process-start
overhead.

\emph{Deeply nested requirements.} With the manual setting ours is faster than
CASAAL on five levels (0.2~s against 0.3~s) and on six (1.3~s against 1.9~s),
and slower on seven (15.1~s against 11.3~s, a factor 1.3) and on eight
(243~s against 73.8~s, a factor 3.3, with 0.43~GB of memory at peak). In
automatic mode ours needs 0.8~s, 6.6~s, 22.1~s and 248~s: the difference to the
manual times is the 5~s budget spent on level \texttt{low}, which is the reason
for the override recommended in Section~\ref{sec:spot-level}, and it vanishes on
N8, where it is below the run-to-run variation. On N7 about 10~s of the 15~s
are Spot's (from a separate run of Spot alone on the simplified formula), which
translates the clocked formula; for six levels this formula already has 99{,}643
characters, and without the proved simplification it is about five times
larger. The reset completion, the optimizer, the export and the initial-value
check take 0.7~s, 2.3~s, 0.6~s and 1.1~s on N7. Two changes on the Spot side
produced these figures (Section~\ref{sec:speedups}): Spot's simulation
reductions are switched off at level \texttt{any}, and the tool reads Spot's
symbolic HOA output instead of cube-expanded LBTT; without them the manual times
would have been 2.2~s, 26.9~s and 1{,}402~s on N6, N7 and N8.

\emph{One driver for both tools.} The figures above compare single runs of ours with the best of three of CASAAL's. We therefore timed both tools with one driver, five runs each, ours in automatic mode under WSL~2 and CASAAL's Windows executable launched from the same shell, on an otherwise idle machine (script \texttt{evaluation/baselines/time\_same\_driver.py}). The medians are 0.85~s (ours) against 0.44~s (CASAAL) on N5, 6.7~s against 2.6~s on N6 and 31~s against 17~s on N7; on the small patterns and the bounded-response conjunctions the two are within about 0.1~s, with ours faster on R5--R8 (for instance 0.07~s against 0.21~s on R7), and a trivial CASAAL run costs 0.15~s of process start. In automatic mode ours is therefore slower than CASAAL on the nested requirements from five levels, and the manual setting narrows this only on five and six levels. The two tools still run on different operating-system kernels.

\emph{Where the time goes on deep nests.} On N7 about 10 of the 15~s are Spot's, which translates the clocked formula: a bounded eventuality repeats the translation of its operand three times, so the formula has size $3^k$ in the nesting depth $k$, and Spot's time grows by a factor of about 16 per level while its output grows by about 3.5. Spot alone takes about as long as CASAAL's whole translation (11.3~s), and the rest of our chain about 5~s, so the difference lies in the replaceable front end (Section~\ref{sec:modular}) and not in the verified optimization and export. The profile, the Spot options we tried and the trade-off behind the automatic default are in Supplement Section~\ref{S-sec:deep-nests}; we did not profile N8.

\subsubsection{Cost for a timed model checker}
\label{sec:casaal-tck}
Supplement Table~\ref{S-tab:casaal-tck} reports the states that TChecker stores
(\texttt{tck-liveness -a couvscc}) on the product of the negated requirement
with a reference system, for the patterns on which both runs completed. The
reference system is the automaton of the requirement itself, which accepts
exactly the words that the negation rejects, so the product is empty and
TChecker explores it completely: the count does not depend on the exploration
order. The same reference is used for every variant. It is exported with the
all-copy policy, because TChecker's liveness algorithms reject diagonal
constraints (it stops with ``unsupported diagonal constraints'') and the
maximal-delay export of the positive form of R2 and the larger bounded-response
patterns contains them (Section~\ref{sec:export-policies}).

Over the 20 completed products TChecker stores 706{,}092 states on our
automata and 1{,}231{,}530 on CASAAL's automata as the adapter encodes them (a
factor 1.7; fewer on 13 patterns, the same on seven, never more). This
comparison mixes two things that the ablations below separate: the translation,
and the way the acceptance condition of CASAAL's automata is given to TChecker.
On the 18 products for which every variant of the ablations completed, the
ratio of CASAAL's count to ours ranges from 1.0 to 8.4 per product, with a
geometric mean of 1.6. The total is dominated by N4, on which both tools store
the same 498{,}087 states; without N4 the factor is 3.5 (208{,}005 against
733{,}443). The largest ratios are on F8 (9 against 76), F7 (280 against 652),
F5 and F6, R3 (168{,}484 against 674{,}698) and R2 (7{,}513 against 21{,}017).
On R3 TChecker needs 12.6~s on our automaton and 23.7~s on CASAAL's, and on N4
149~s and 231~s; these times are single runs on a machine that was also running
other jobs, so we do not read them beyond the direction. On R4--R7 and N5--N6
TChecker did not finish within 300~s for either tool (on R4 it had used
1.7~GB), and for R8, N7 and N8 the reference system or the automaton could not
be generated within the limit, so nothing is claimed for those.

\subsubsection{Where the difference comes from}
\label{sec:casaal-ablation}
\begin{table}[tp]
\centering
\caption{TChecker stored states on the product of the negated requirement with
the reference system (setup of Supplement Table~\ref{S-tab:casaal-tck}), summed over the 18
products on which every variant completed (F1--F12 except the hatted F3,
R1--R3, N1--N3), and over three subsets. The last two rows pass CASAAL's
automaton through a small importer (outside the proofs) that converts its
transition acceptance to location acceptance, and then through our export.
``Invariants removed'' deletes every location invariant from the automaton
given to TChecker; the ``invariant script'' sets, at each location, $x\le k$
for the largest upper bound $k$ of the outgoing guards when every outgoing
guard bounds $x$ (no propagation, no splitting, unverified).}
\label{tab:baseline-ablation}
\small
\begin{tabular}{@{}lrrrr@{}}
\toprule
Automaton given to TChecker & all 18 products & F5--F8 & F10--F12 & R3 \\
\midrule
ours (default export) & 204{,}564 & 771 & 2{,}943 & 168{,}484 \\
ours, invariants removed & 204{,}749 & 771 & 3{,}128 & 168{,}484 \\
CASAAL, transition acceptance encoded by the adapter & 726{,}543 & 1{,}708 & 3{,}835 & 674{,}698 \\
CASAAL + invariant script & 726{,}358 & 1{,}708 & 3{,}650 & 674{,}698 \\
CASAAL converted to location acceptance, our export & 204{,}819 & 1{,}026 & 2{,}943 & 168{,}484 \\
same, invariants removed & 205{,}004 & 1{,}026 & 3{,}128 & 168{,}484 \\
\bottomrule
\end{tabular}

\end{table}
Four ablations separate the stages (Table~\ref{tab:baseline-ablation}).
\emph{Optimizer.} With the optimizer disabled (\texttt{-n 0}, last column of
Supplement Table~\ref{S-tab:casaal-tck}) TChecker stores exactly the same number of states
on each of the 20 products: the optimizer reduces the size of the automaton,
from 688 to 461 transitions on N5, from 2{,}737 to 1{,}541 on N6, from 10{,}930
to 5{,}105 on N7 (the size of CASAAL's automata in each case) and, on the
positive form of five conjoined bounded responses, the exported automaton from
5{,}469 to 651 transitions, but it does not reduce what TChecker explores on
these products. Running TChecker on the automaton before any pass (the output
of reset completion) gives the same count as on ours on six of seven products
tried (F1, F7, R2, R3, N2, N3) and 1{,}790 instead of 1{,}688 on F12.

\emph{Invariants.} Only F10--F12 carry location invariants in our negated
automata; the others have none. Removing the invariants of our own export
changes the total from 204{,}564 to 204{,}749, all of it on F10--F12. Giving
CASAAL's automata the invariant script changes 726{,}543 to 726{,}358. The
location invariants of the export therefore explain at most 185 of the
204{,}564 states, and they are not the source of the factor 1.7.

\emph{Acceptance encoding.} TChecker has labels on locations only, so CASAAL's
transition-based acceptance has to be encoded. The adapter inserts a committed
location after every accepting transition, and each such location is a stored
state of the product. Converting the acceptance to locations, as the importer
does, and exporting with our tool gives 204{,}819 states for the 18 products
instead of 726{,}543, against 204{,}564 for our translation. The gap closes on
all but four patterns; the 255 states that remain are on F5--F8, where our
translation gives a smaller automaton than CASAAL's (on F7 484 against 280 and
on F8 54 against 9). The importer reads CASAAL's DOT output and, when CASAAL
prints no acceptance mark, treats a state as accepting unless its label
contains an unbounded eventuality, as the adapter does. A first version of the
importer ignored that convention and gave an empty accepting set on F9--F11 and
N1--N4; the verdicts of the non-empty workload below exposed this, and the
state counts of the empty products were unaffected. The importer did not
convert F14, which has three acceptance sets, and the driver did not finish on
N6--N8 within 300~s.

\emph{Optimizer on CASAAL's automata.} On every CASAAL automaton that the
importer converts (F1--F12, R1--R8, N1--N5), the optimizer changes nothing:
these automata are a fixed point of our passes, which is consistent with their
equal size. A hand-made automaton with redundant guards shrinks under the same
driver, so the importer does not prevent the passes from acting.

The advantage over CASAAL as encoded is therefore a property of the acceptance
representation, not of the invariants and not of the optimizer. What remains of
our translation is the 255 states on F5--F8 and the smaller automata on F8,
F11, F12 and F14; the verified export gives the same TChecker counts as an
unverified invariant script, and its assurance comes at no measurable cost.

\subsubsection{Non-empty products}
\label{sec:casaal-nonempty}
All products above are empty, so they show the size of the state space but not
that the translators agree on a verdict when the system violates the
requirement. We therefore composed the negated automata with two kinds of
system. A \emph{free} system allows every event at every time (strictly
increasing, divergent timestamps): for each of the 20 products, ours (default,
\texttt{-n 0}, invariants removed), CASAAL's own automaton (with and without
the invariant script) and CASAAL converted and exported give a non-empty
product, 138 verdicts in all, with no disagreement (F14 cannot be converted).
A \emph{response} system for R1--R3 either answers every request within two
time units or never answers the first one; the products are empty in the first
case and non-empty in the second, for every variant and for both export
policies, with no disagreement (42 verdicts for each case). TChecker stops at
the first violation, so state counts of non-empty products depend on the search
order and are not compared. Section~\ref{sec:verification-adds} reports the same check
with UPPAAL.

\subsubsection{Summary of the comparison}
The automata have the same size on every bounded-response and nested pattern both tools complete, and ours are smaller on three small ones. TChecker stores fewer states on ours than on CASAAL's automata as encoded for it (1.7 times fewer in total), but almost all of this disappears once CASAAL's acceptance is converted to locations; our invariants account for at most 185 states and the optimizer for none. In generation time ours is faster on the bounded-response conjunctions and slower on nested requirements, where Spot is the bottleneck.
